\section{A symmetric algebra and a degree-three class}
\label{alg:section}

We begin with a symmetric algebra carrying a simple module and a
Hochschild class of degree three.  The class will have two uses: its
value on the simple will generate the positive self-extension algebra,
and an explicit boundary identity will allow us to lift maps through
the two cones used later.

\subsection{The ten-dimensional algebra and its trivial extension}

Let $C$ be the vector space with the following basis, arranged by
corners.  The notation $iCj$ specifies the left idempotent first.
\begin{equation}
\begin{array}{c|c@{\qquad}c|c}
 \text{corner}&\text{basis}&\text{corner}&\text{basis}\\ \hline
 eCe&e,x,y,z&eCf&u,v\\
 fCe&t,j&fCf&f,n
\end{array}
\label{alg:C-corners}
\end{equation}
The elements $e,f$ are orthogonal idempotents with sum $1$.  Besides
their actions specified by the corners, the complete list of nonzero
products is
\begin{equation}
\begin{gathered}
 xy=qz,\qquad yx=z,\qquad xu=yu=v,\qquad tx=j,\qquad ty=q^2j,\\
 ut=y+qx,\qquad vt=qz,\qquad uj=z,\qquad
 tu=(1+q)n,\qquad nt=qj,\qquad un=v.
\end{gathered}
\label{alg:C-products}
\end{equation}
All products not specified by these rules are zero.

\begin{lemma}\label{alg:C}
The rules \eqref{alg:C-corners}--\eqref{alg:C-products} define an
associative unital algebra of dimension ten.  Its radical is the span
of the eight non-idempotent basis vectors, and
\[
 (\rad C)^5=0,\qquad utut=q(1+q)z\ne0.
\]
\end{lemma}

\begin{proof}
Give $u,t$ degree one, $x,y,n$ degree two, $v,j$ degree three, and $z$
degree four; the idempotents have degree zero.  Every listed product
is homogeneous.  Associativity involving an idempotent follows from
the corners.  For three non-idempotent basis vectors, incompatible
corners or total degree greater than four force both bracketings to
vanish.  The complete remaining list is recorded below; direct
substitution gives the same value under either bracketing.
\[
\begin{array}{c|ccccccccccc}
 abc&utu&tut&xut&yut&ntu&unt&txu&tyu&utx&uty&tun\\ \hline
 (ab)c=a(bc)&(1+q)v&q(1+q)j&qz&qz&0&qz&0&0&z&q^2z&0
\end{array}
\]
This proves associativity by trilinearity.  The positive-degree span
$N$ is a nilpotent ideal with $N^5=0$ and $C/N\cong\kk\times\kk$;
therefore $N=\rad C$.  Finally, $(ut)u=(1+q)v$ and $vt=qz$ give the
asserted nonzero fourth radical product.
\end{proof}

\begin{remark}\label{alg:antecedent}
The corner $eCe$ is the quantum exterior algebra
$\kk\langle x,y\rangle/(x^2,y^2,xy-qyx)$, with $z=yx$.
The right module $fCe$ has basis $t,j$ and actions
$tx=j$, $ty=q^2j$, $jx=jy=0$. More precisely, put $R=eCe$.
The map
\[
 t\longmapsto x+q^{-1}y,\qquad j\longmapsto q^{-1}z
\]
identifies $fCe$ with the right ideal $(x+q^{-1}y)R$.
In the notation of Schulz \cite[Example~7]{Schulz1995}, the parameter
is $q^{-1}$ and this ideal is the first syzygy of $(x+y)R$.
Here its interactions with the other corners are fixed by the complete
multiplication table above.
\end{remark}

Put $L=DC$ and equip it with the dual bimodule actions
\begin{equation}
 (a\phi)(c)=\phi(ca),\qquad (\phi a)(c)=\phi(ac)
 \quad(a,c\in C,\ \phi\in L).
\label{alg:dual-actions}
\end{equation}
Trivial extensions and their cohomology are studied in \cite{CibilsMarcosRedondoSolotar2003}. Here the trivial extension is
\begin{equation}
 T=C\ltimes L,
 \qquad (a,\phi)(b,\psi)=(ab,a\psi+\phi b).
\label{alg:T-definition}
\end{equation}
For every lower-case basis letter $a$ of $C$, its upper-case letter
$a^*$ denotes the dual basis vector in $L$; for example, $j^*=J$ and
$f^*=F$.  We also set $(a^*)^*=a$.  Thus a capital letter $J$ is a
basis vector, whereas the radical below is denoted by $\mathfrak r$.

\begin{lemma}\label{alg:T}
The algebra $T$ is symmetric of dimension twenty.  A symmetrizing
trace is
\[
 \lambda(a,\phi)=\phi(1),
\]
and the trace-dual of each basis letter is its starred letter.  Its
radical $\mathfrak r$ is the span of all basis letters except $e,f$;
it satisfies $\mathfrak r^6=0$ and $T/\mathfrak r\cong\kk\times\kk$.
\end{lemma}

\begin{proof}
The dual actions satisfy the bimodule identities by associativity of
$C$, so \eqref{alg:T-definition} is associative.  Moreover,
\[
 \lambda\bigl((a,\phi)(b,\psi)\bigr)=\psi(a)+\phi(b),
\]
which is symmetric and nondegenerate.  In particular,
$\lambda(ab^*)=\delta_{a,b}$ for lower-case basis letters $a,b$.
The dual of the corner $iCj$ lies in $jLi$.

Extend the grading of $C$ by
$\deg(a^*)=5-\deg(a)$.  Formula \eqref{alg:dual-actions} shows that
every product in $T$ is homogeneous.  All basis vectors except $e,f$
have positive degree, and no degree exceeds five.  Their span is
therefore a nilpotent ideal whose quotient is $\kk\times\kk$;
this proves the radical assertions.
\end{proof}

Let $s$ be the one-dimensional left $T$-module on which $f$ acts as
$1$ and every other basis letter acts as zero.  It is inflated from
the corresponding $C$-simple.  Write $s^{\mathrm r}$ for the right
simple defined by the same character.

\subsection{The bar class and the finite identities}

Set $I=\kk e\oplus\kk f$.  The normalized relative bar resolution of
$T$ has terms
\begin{equation}
 P_n=T\otimes_I\mathfrak r^{\otimes_I n}\otimes_I T
 \quad(n\ge0),\qquad P_n=0\quad(n<0).
\label{alg:bar}
\end{equation}
We write $a_0[a_1|\cdots|a_n]a_{n+1}$ for an elementary tensor.  The
differential multiplies each adjacent pair of factors and sums the
results; the augmentation $P_0\to T$ is multiplication.  All signs
are $+$ because the characteristic is two.  Splitting by idempotents
expresses each term as a finite sum of projective bimodules
$Ti\otimes jT$ with $i,j\in\{e,f\}$.  Exactness follows from a contraction:
replace the first coefficient $a_0$ by its projection to
$\mathfrak r$, insert that projection as a new first bar, and put
$1$ in the first coefficient.  In augmented degree send $a$ to
$1[]a$.  Writing $\sigma$ for this contraction, the identity
$d\sigma+\sigma d=\id$ follows by canceling successive
multiplications; the $I$-component of $a_0$ is balanced over $I$.
This contraction is right $T$-linear, so $P\otimes_T s$ is a
projective resolution of $s$ as well.

Give a lower-case letter weight zero and an upper-case letter weight
one.  This weight records the effect of scaling the square-zero ideal
$L=DC$; it is distinct from the positive grading in
Lemma~\ref{alg:T}.  The complete table in Appendix~\ref{app:cochain}
defines an $I$-bimodule map
\[
 p:\mathfrak r\otimes_I\mathfrak r\otimes_I\mathfrak r\longrightarrow T.
\]
Every unlisted value is zero.  The following records exactly the
properties of this finite table that the proof needs.

\begin{lemma}\label{coc:data}
The cochain $p$ has weight $-1$ and is a Hochschild cocycle.  Its
associated class evaluated on $s$ is nonzero in $\Ext_T^3(s,s)$.

For any $H\in\kk^\times$, let $h=h_H$ be the automorphism of $T$
fixing $C$ and multiplying $L$ by $H$.  On the radical basis define
\[
 z_H(a)=
 \begin{cases}
 Hqa,&a\in\{u,v,t,E,X,Y,Z,U,V\},\\
 0,&\text{otherwise},
 \end{cases}
\]
and extend linearly.  For all $a,b\in\mathfrak r$ one has
\begin{equation}
 \sum_{w\in\mathcal B_{\mathfrak r}}
      p(a,b,h(w))w^*
  =a z_H(b)+z_H(ab)+z_H(a)h^{-1}(b),
\label{coc:boundary}
\end{equation}
where $\mathcal B_{\mathfrak r}$ is the radical basis.
\end{lemma}

\begin{proof}
The endpoint, weight, cocycle and boundary identities are verified
from the complete finite data in Appendix~\ref{app:cochain}.
We prove the nonzero evaluation explicitly.  In the bar complex
$s^{\mathrm r}\otimes_T P\otimes_T s$, consider
\begin{equation}
 \zeta=q^2[t|x|J]+[t|y|J].
\label{coc:cycle}
\end{equation}
Both words have endpoint $f$ on each side.  All outer radical
actions vanish on the endpoint simples.  The multiplication rules,
including the dual actions, give
\[
 tx=j,\qquad ty=q^2j,\qquad xJ=T,\qquad yJ=q^2T.
\]
Consequently
\[
 d\zeta=q^2([j|J]+[t|T])+q^2[j|J]+q^2[t|T]=0.
\]
The table gives $p(t,x,J)=0$ and $p(t,y,J)=q^3f$, so its evaluation on
$\zeta$ is $q^3\ne0$.  Every $T$-linear map from $P\otimes_T s$ to
$s$ factors through the left simple quotient and hence pairs with
this two-sided-simple bar complex.  A coboundary pairs to zero with
its cycles.  Thus the evaluated class is nonzero.
\end{proof}

The evaluation of a Hochschild class on modules is the characteristic action; see \cite[Section 4.3]{Lowen2008}. We use the same letter $p$ for its comparison map $P\to P[3]$:
for $n\ge3$ apply the cochain to the last three bars and multiply its
output into the right coefficient,
\begin{equation}
 a_0[a_1|\cdots|a_n]a_{n+1}\longmapsto
 a_0[a_1|\cdots|a_{n-3}]
 p(a_{n-2},a_{n-1},a_n)a_{n+1};
\label{coc:comparison}
\end{equation}
the map is zero for $n<3$.  This is a chain map: strict-prefix
multiplications cancel, and the remaining equation on the last
four bars is precisely the Hochschild cocycle equation.  At the
bottom index the target differential is zero.  Composing its
evaluation on $s$ with the augmentation represents the class in
Lemma~\ref{coc:data}.

We have now obtained the finite algebra and the specific degree-three
class.  The next step is to compute all self-extension degrees of
$s$, so that this nonzero class can be identified as a polynomial
generator.
